\subsection{\aijon}
\label{aijon}

\aijon is a fuzzer that implements IJON\cite{aschermann2020ijon} annotations on top of AFL++.
It contains an LLM agent which inserts these annotations into the target's source.

\aijon takes as input a list of functions that are identified by \codeswipe to possibly contain a vulnerability.
It then uses the LLM agent to analyze the source code of this single function, so as to reduce the LLM context, and inserts IJON style annotations into its source code.

The LLM is instructed to strategically select annotations that can guide the fuzzer towards the vulnerability specified in the \codeswipe report.
These annotations provide valuable feedback to AFL++ that can guide it very quickly to the target locations and find a crashing input if any for the specified vulnerability.

However, using such annotations can lead to collisions within the bitmap of AFL++.
In order to address this problem, we optimize the instrumentation step of AFL++ by limiting the instrumentation to a select list of functions.
\aijon uses the callgraph contained in the \analysisgraph to identify the functions on the path from the harness' entry point to the target function.
It then instructs AFL++ to only instrument the functions that are present in these paths.
This allows AFL++ to focus on the vulnerability in the target function without exploring other sections of the target program and reduces the possibility of collisions in the bitmap.

\aijon also uses the \analysisgraph to generate a starting seed corpus for fuzzing.
Other fuzzers, such as \afluzzer, could have accidentally generated inputs that reach some of the functions in the call path to the target function.
These inputs are optimal candidates for being used as seeds, since they have triggered a function close to the target function in the callgraph.
\aijon collects such inputs from \analysisgraph and uses them as the starting seed corpus.

Any interesting inputs or crashes found by \aijon are synced with \afluzzer.
This allows \afluzzer to replicate the crashes on the target without annotations and acts as a verification step for the crashes detected by \aijon.